%document
\documentclass[10pt]{beamer}
%theme
\usetheme{metropolis}
% packages
\usepackage{color}
\usepackage{listings}
\usepackage[ngerman]{babel}
\usepackage[utf8x]{inputenc}
\usepackage{multicol}


% color definitions
\definecolor{mygreen}{rgb}{0,0.6,0}
\definecolor{mygray}{rgb}{0.5,0.5,0.5}
\definecolor{mymauve}{rgb}{0.58,0,0.82}

\lstset{
    backgroundcolor=\color{white},
    % choose the background color;
    % you must add \usepackage{color} or \usepackage{xcolor}
    basicstyle=\footnotesize\ttfamily,
    % the size of the fonts that are used for the code
    breakatwhitespace=false,
    % sets if automatic breaks should only happen at whitespace
    breaklines=true,                 % sets automatic line breaking
    captionpos=b,                    % sets the caption-position to bottom
    commentstyle=\color{mygreen},    % comment style
    % deletekeywords={...},
    % if you want to delete keywords from the given language
    extendedchars=true,
    % lets you use non-ASCII characters;
    % for 8-bits encodings only, does not work with UTF-8
    frame=single,                    % adds a frame around the code
    keepspaces=true,
    % keeps spaces in text,
    % useful for keeping indentation of code
    % (possibly needs columns=flexible)
    keywordstyle=\color{blue},       % keyword style
    % morekeywords={*,...},
    % if you want to add more keywords to the set
    numbers=left,
    % where to put the line-numbers; possible values are (none, left, right)
    numbersep=5pt,
    % how far the line-numbers are from the code
    numberstyle=\tiny\color{mygray},
    % the style that is used for the line-numbers
    rulecolor=\color{black},
    % if not set, the frame-color may be changed on line-breaks
    % within not-black text (e.g. comments (green here))
    stepnumber=1,
    % the step between two line-numbers.
    % If it's 1, each line will be numbered
    stringstyle=\color{mymauve},     % string literal style
    tabsize=4,                       % sets default tabsize to 4 spaces
    % show the filename of files included with \lstinputlisting;
    % also try caption instead of title
    language = Java,
	showspaces = false,
	showtabs = false,
	showstringspaces = false,
	escapechar = ,
}

\def\ContinueLineNumber{\lstset{firstnumber=last}}
\def\StartLineAt#1{\lstset{firstnumber=#1}}
\let\numberLineAt\StartLineAt



\newcommand{\codeline}[1]{
	\alert{\texttt{#1}}
}

% This Document contains the information about this course.

% Authors of the slides
\author{Fritz Windisch, Lukas Tassanyi}

% Name of the Course
\institute{Java-Course}

% Fancy Logo
\titlegraphic{\hfill\includegraphics[height=1.25cm]{../templates/fsr_logo_cropped}}


\title{Java}
\subtitle{Introduction}
\date{\today}


\begin{document}

\begin{frame}
    \titlepage
\end{frame}
\begin{frame}{Overview}
    \setbeamertemplate{section in toc}[sections numbered]
    \tableofcontents
\end{frame}

\section{Proceeding}
\begin{frame}{About this course}
    Requirements
    \begin{itemize}
               \item Please bring your computer with You or become familiar with the PC Pool
        \item You know how to use a computer
        \item Basic knowledge in mathematics and logic
                \item Installed JDK and IDE (BlueJ/IntelliJ/Eclipse)
        \item Maybe already knowledge in programming languages?
    \end{itemize}
    Proceeding
    \begin{itemize}
        \item There will be ~14 lessons
        \item Each covers a topic and comes with excercises
    \end{itemize}
\end{frame}

\begin{frame}{Some resources}
    \begin{itemize}
        \item You can ask your tutor
        \item Join the Auditorium group \hfill \\
            \url{https://auditorium.inf.tu-dresden.de/en/groups/110804211}
        \item StackOverflow, FAQs, Online-tutorials, ... \hfill \\
        \item Official documentation \hfill \\
            \url{https://docs.oracle.com/javase/13/}
           \item mailinglist \url{programmierung@ifsr.de}
    \item Material-Repository \\
        \url{https://friwi.me/java-lessons/}
    \end{itemize}
\end{frame}

\begin{frame}{About Java}
    Pros:
    \begin{itemize}
        \item Syntax like C++
        \item Strongly encourages OOP
        \item Platform-independent (JVM)
        \item Very few external libraries
        \item[] $->$ Easy to use and very little to worry about
    \end{itemize}
\end{frame}

\begin{frame}{About Java}
    Cons:
    \begin{itemize}
        \item A lot of unnecessary features \\ in the JDK
        \item Slower than assembly
        \item No multi-inheritance
        \item Mediocre support for other programming paradigms
        \item[] $->$ Neither fast, small nor geeky
    \end{itemize}
\end{frame}

\section{Your first program}
\begin{frame}[fragile]{Creating your Working Environment}
  \begin{itemize}
      \item Open BlueJ and create a new project\\
          \item IntelliJ users will need to take different steps (just try on your own or ask us)
      \item Create a new class, name it ``Hello''
      \item Right click on your new class, choose ``Edit''
  \end{itemize}
\end{frame}

\subsection{Hello World!}

\begin{frame}[fragile]{Hello World!}
    This is an empty Java class.
      Java classes always start with a capital letter.
      Remove the initial contents of the file and replace it with the following code:
    \begin{lstlisting}
    public class Hello {

    }\end{lstlisting}
\end{frame}

\begin{frame}[fragile]{Hello World!}
    This is a small program printing \emph{Hello World!} to the console:
    \begin{lstlisting}
    public class Hello {
        public static void main(String[] args) {
            System.out.println("Hello World!");
        }
    }\end{lstlisting}
\end{frame}

\begin{frame}[fragile]{How to run your program}
   \begin{itemize}
      \item Click on ``translate'' in the class edit dialogue
      \item Check if there have been any compile errors
      \item If not, right-click the class in the root window
      \item Choose the ``void main(String[] args)'' method to be executed
      \item In the appearing dialogue, just press ``OK''
  \end{itemize}
\end{frame}

\section{Basics}

\begin{frame}[fragile]{Comments}
   \begin{lstlisting}
   public class Hello {
       // a single line comment
       public static void main(String[] args) {
           System.out.println("Hello World!");
           /* a comment
           spanning
           multiple lines
           */
       }
   }\end{lstlisting}
   \begin{itemize}
       \item // endet am Zeilenende
       \item /* endet mit */
   \end{itemize}
\end{frame}

\subsection{Some definitions}

\begin{frame}[fragile]{Code concepts}
    \begin{lstlisting}
    public class Hello {
        public static void main(String[] args) {            
            int x;
            x = 9;
            int y = 23;
            int z;
            z = x * y;
            
            System.out.println(z);
        }
    }\end{lstlisting}
\end{frame}

\begin{frame}[fragile]{Code concepts}
    \begin{lstlisting}
    public class Hello {
        public static void main(String[] args) {
            System.out.println(9 * 23);
        }
    }\end{lstlisting}
\end{frame}

\begin{frame}{Primitive data types}
    Java supports some primitive data types:
    \begin{itemize}
        \item[boolean] a truth value (either \textbf{true} or \textbf{false})
        \item[int] a 32 bit integer (e.g. 127 or -2147483648)
        \item[long] a 64 bit integer (e.g. -9,223,372,036,854,775,808)
        \item[float] a 32 bit floating point number (e.g. 3.14159F)
        \item[double] a 64 bit floating point number (e.g. 3.15159)
        \item[char] an ascii character (e.g. \textquotesingle A\textquotesingle, \textquotesingle \_\textquotesingle, \textquotesingle \enspace\textquotesingle)
        \item[void] the empty type (needed in later topics)
    \end{itemize}
\end{frame}

\begin{frame}[fragile]{About the Semicolon}
    \begin{lstlisting}
    public class Hello {
        // prints a "Hello World!" on your console
        public static void main(String[] args) {
            System.out.println("Hello World!");
        }
    }\end{lstlisting}
    Semicolons conclude all statements. \\
    Blocks do not need a semicolon.
\end{frame}

\begin{frame}[fragile]{Blocks}
    \begin{lstlisting}
    public class Hello {
    
        public static void main(String[] args) {
            System.out.println("Hello World!");
        }
        
    }\end{lstlisting}
    Everything between \{ and \} is a \emph{block}. \\
    Blocks may be nested.
\end{frame}

\begin{frame}[fragile]{Naming of Variables}
    \begin{itemize}
        \item The names of variables can begin with any letter or underscore. \\
        \item Variables and functions start with small, classes with capital letters.
        \item Compound names should use CamelCase.
        \item Use meaningful names.
    \end{itemize}
    \begin{lstlisting}
    public class Calc {
        public static void main(String[] args) {
            int a = 0; // not very meaningful
            float myFloat = 5.3f; // also not meaningfull
            int count = 7; // quite a good name

            int rotationCount = 7; // there you go
        }
    }\end{lstlisting}
\end{frame}

\subsection{Calculating}

\begin{frame}[fragile, allowframebreaks]{Calculating with \emph{int}}
    \begin{lstlisting}
    public class Calc {
        public static void main(String[] args) {
            int a; // declare variable a
            a = 7; // assign 7 to variable a
            System.out.println(a); // prints: 7
            a = 8;
            System.out.println(a); // prints: 8
            a = a + 2;
            System.out.println(a); // prints: 10
        }
    }\end{lstlisting}
    After the first assignment the variable is initialized.
\framebreak
    \begin{lstlisting}
    public class Calc {
        public static void main(String[] args) {
            int a = -9; // declaration and assignment of a
            int b; // declaration of b
            b = a; // assignment of b
            System.out.println(a); // prints: -9
            System.out.println(b); // prints: -9
            a++; // increments a
            System.out.println(a); // prints: -8
        }
    }\end{lstlisting}
\framebreak
    Some basic mathematical operations:
    \begin{tabular}{ll}
        Addition & \texttt{a + b} \\
        Subtraction & \texttt{a - b} \\
        Multiplication &\texttt{a * b} \\
        Division & \texttt{a / b} \\
        Modulo & \texttt{a \% b} \\
        Increment & \texttt{a++;} \\
        Decrement & \texttt{a--;} \\
    \end{tabular}
\end{frame}

\begin{frame}[fragile, allowframebreaks]{Calculating with \emph{float}}
    \begin{lstlisting}
    public class Calc {
        public static void main(String[] args) {
            float a = 9;
            float b = 7.5f;
            System.out.println(a); // prints: 9.0
            System.out.println(b); // prints: 7.5
            System.out.println(a + b); // prints: 16.5
        }
    }\end{lstlisting}
\framebreak
    \begin{lstlisting}
    public class Calc {
        public static void main(String[] args) {
            float a =       8.9f;
            float b = 3054062.5f;
            System.out.println(a); // prints: 8.9
            System.out.println(b); // prints: 3054062.5
            System.out.println(a + b); // prints: 3054071.5
        }
    }\end{lstlisting}
    Float has a limited precision. \\
    \emph{This might lead to unexpected results!}
\end{frame}

\begin{frame}[fragile]{Mixing \emph{int} and \emph{float}}
    \begin{lstlisting}
    public class Calc {
        public static void main(String[] args) {
            float a = 9.3f;
            int b = 3;
            System.out.println(a + b); // prints: 12.3
            float c = a + b;
            System.out.println(c); // prints: 12.3
        }
    }\end{lstlisting}
    Java converts from \textbf{int} to \textbf{float} by default, if necessary. \\
    But not vice versa.
\end{frame}

\subsection{Text with Strings}

\begin{frame}[fragile]{Strings}
    A String is not a primitive data type, but an object. \\
    We discuss objects in detail in the next lesson.
    \begin{lstlisting}
    public class Calc {
        public static void main(String[] args) {
            String hello = "Hello World!";
            System.out.println(hello); // print: Hello World!
        }
    }\end{lstlisting}
\end{frame}

\begin{frame}[fragile]{Concatenation}
    \begin{lstlisting}
    public class Calc {
        public static void main(String[] args) {
            String hello = "Hello";
            String world = " World!";
            String sentence = hello + world;
            System.out.println(sentence);
            System.out.println(hello + " World!");
        }
    }\end{lstlisting}
    You can concatenate Strings using the +. Both printed lines look the same.
\end{frame}

\begin{frame}[fragile]{Strings and Numbers}
    \begin{lstlisting}
    public class Calc {
        public static void main(String[] args) {
            int factorA = 3;
            int factorB = 7;
            int product = factorA * factorB;
            String answer =
                factorA + " * " + factorB + " = " + product;
            System.out.println(answer); // prints: 3 * 7 = 21
        }
    }\end{lstlisting}
    Upon concatenation, primitive types will be replaced by their current value as \emph{String}.
\end{frame}

\begin{frame}[fragile]{Exercise and questions}
    \begin{itemize}
    \item You can now begin solving the first exercise from the course materials yourself or with a partner. \\
        \url{https://friwi.me/java-lessons/}
    \item If you require any assistance, feel free to ask us at any time
    \end{itemize}
\end{frame}

\begin{frame}[fragile]{Thank you for your attention}
    \begin{itemize}
    \item Course materials: \\
        \url{https://friwi.me/java-lessons/}
    \end{itemize}
\end{frame}

\end{document}

